% ============================================================
% M3 S3 练习件·W2 臂B —— 课时08 直线与圆的位置关系（2.3.3）＝照抄母版件（课时01）体例
% 生成：M3 成卷轮 S3 W2 臂B｜2026-09-14｜编译 cwd＝本目录（中文路径禁 \input 绝对路径）
%   xelatex main-true.tex ×2 → 含详解印本；xelatex main-false.tex ×2 → 纯题版（单源双本）
% 输入链（全只读）：题面库 课时08-直线与圆的位置关系.md（题面逐字装配）＋ -答案侧.md（值逐字回嵌）
%   ＋ 定稿/批B-课时08-直线与圆的位置关系.md（详解＝件内【亲算】格回嵌源）；toolchain＝qp-m3.sty
%   （钉后版 md5 7c3930362be8a0a2bdf21bbf8ac16573，本地挂载副本，破损即停）。
% 母版体例：详见《练习件母版体例说明.md》（课时01 件）；门谱读数见 过程件 工作区/_tmpM3S3W2臂B0914/。
% 键账：件内 ansblock 键集 ≡ 题面库片练键集 ≡ manifest 练键序 T1~T16（16 键，零缺漏零多余）；
%   拓区隔离：2章-拓/测/滚 键不入本件（S4 拓展册收）；导学键 G1~G5 属导学件域，亦不入。
% 卷式例外案B：练习件不属卷式——答案块物理落点恒为题后 ansblock，禁卷末附卷制；
%   全线括线模（导言 \ansblockgrayfalse 一次置定，件内不切换；灰底盒不可跨栏断，练习件禁用）。
% 题序＝manifest 键序 T1~T16（零重排）；印面题号 1~16 连号（\tihao 号＝\ansitem 号同键）；
%   三档切位 1~7／8~14／15~16（照库序切位，非难度重排）。
% 槽配实录：单选10（T1~T10）＋多选0＋填空4（T11~T14）＋解答2（T15/T16）；
%   10选4填2解 为波次方案基线槽配，本片实配与基线一致（批B §四 满足度表同口径）。
% 选项槽位：默认四连排（0.25\linewidth−0.25\lxhang），超宽降档梯 四连排→两连排
%   （0.5\linewidth−0.5\lxhang−0.25em）→单列 \lxopt；禁缩字号/负 kern/删标点凑宽。
%   本片实测降档：见值台账「降档登记」＋回执（槽宽门真算读数）。
% 解答书写区：T15 三问 32mm；T16 两问 24mm（16~40mm 带内；纯题档吞 ansblock 不吞 \liubai）。
% 填空印答：T11~T14 空位 \kongda{值}（详解版印答、纯题版退定宽空线，两档等形）。
% 难度词：题面侧头标第 4 字段系裸数值（批B 口径），派生标注照同章导学件09 同族口径
%   0.94/0.85→简单、0.65→中档（派生登记见值台账）；知识点 N 照导学件08 知识导学三分册：
%   一＝直线与圆的位置关系、二＝圆的弦长问题、三＝圆的切线方程。
% 印面清洗：T15 亲算行 ✓ 记号剔（检验句文句保留）；题面侧 T15/T16【注】行不入印面；
%   T16 题面去示意图引用（源含图，全部数学信息在文字中，批B 台账同注）。
% ============================================================
\documentclass[fontset=none]{ctexart}
\usepackage{qp-m3}
% —— 印面逐字符号路由补钉（承导学件母版；− 走 TNR、▱ NSC 子块；禁改 sty 保 md5） ——
\xeCJKDeclareCharClass{Default}{"2212}%
\setCJKfamilyfont{nscblk}[Path=C:/提示词/工作区/字替对照-0909/variantF/fonts/,BoldFont=NSC-w600.ttf]{NSC-w500.ttf}%
\xeCJKDeclareSubCJKBlock{nscblk}{"25B1}%
\newcommand{\pxparallelogram}{{\CJKfamily{nscblk}▱}}%
% —— 断行微调（承导学件母版实测档，三零门承重；\emergencystretch 不另设） ——
\xeCJKsetup{CJKglue={\hskip 0.10em plus 0.08em minus 0.05em}}%
% —— 练习件全线括线模（S3 波次方案 §四.3：导言一次置定、件内不切换） ——
\ansblockgrayfalse
% —— 页脚件名（qp-headfoot 复用入口） ——
\renewcommand{\qpzhangming}{第二章\quad 平面解析几何}%
\renewcommand{\qpjianming}{练习件}%

\begin{document}

\keshi{第8课时\quad 直线与圆的位置关系}
\begin{multicols}{2}
\emergencystretch=1em

% ============================================================
% 基础巩固（题1~7＝T1~T7：单选7；题后紧跟答案制，全线括线 ansblock）
% ============================================================
\dangbadge{基}{础}{巩}{固}

\tihao{1}\zu{圆的弦长问题×单选} \tieside{简单(知识点二)}直线\(l\)：\(3x+4y-1=0\)被圆\(C\)：\(x^{2}+y^{2}-2x-4y-4=0\)所截得的弦长为\nobreak(\kongwei)
\bindopt {\fontsize{10.5pt}{\qpTJgLeadLiti}\selectfont \makebox[\dimexpr0.25\linewidth-0.25\lxhang\relax][l]{A．\(2\sqrt{5}\)}\makebox[\dimexpr0.25\linewidth-0.25\lxhang\relax][l]{B．\(4\)}\makebox[\dimexpr0.25\linewidth-0.25\lxhang\relax][l]{C．\(2\sqrt{3}\)}\makebox[\dimexpr0.25\linewidth-0.25\lxhang\relax][l]{D．\(2\sqrt{2}\)}\par}

\begin{ansblock}[2章-练-课时08-T1]
% ans:2章-练-课时08-T1
\ansitem{1}{A}
\ansnote{详解}{圆\(C\)即\((x-1)^{2}+(y-2)^{2}=9\)，圆心\((1,2)\)，\(r=3\)．圆心到直线\(l\)的距离\(d=|3+8-1|/5=2\)．由弦长公式，弦长\(=2\sqrt{9-4}=2\sqrt{5}\)．故选：A．}
\end{ansblock}

\tihao{2}\zu{圆的弦长问题×单选} \tieside{中档(知识点二)}已知圆\(x^{2}+y^{2}=9\)的弦过点\(P(1,2)\)，当弦长最短时，该弦所在直线方程为\nobreak(\kongwei)
\bindopt {\fontsize{10.5pt}{\qpTJgLeadLiti}\selectfont \makebox[\dimexpr0.5\linewidth-0.5\lxhang-0.25em\relax][l]{A．\(x+2y-5=0\)}\makebox[\dimexpr0.5\linewidth-0.5\lxhang-0.25em\relax][l]{B．\(y-2=0\)}\par}
\bindopt {\fontsize{10.5pt}{\qpTJgLeadLiti}\selectfont \makebox[\dimexpr0.5\linewidth-0.5\lxhang-0.25em\relax][l]{C．\(2x-y=0\)}\makebox[\dimexpr0.5\linewidth-0.5\lxhang-0.25em\relax][l]{D．\(x-1=0\)}\par}

\begin{ansblock}[2章-练-课时08-T2]
% ans:2章-练-课时08-T2
\ansitem{2}{A}
\ansnote{详解}{过圆内点的弦以垂直于\(OP\)的弦最短．\(k_{OP}=2\)，故所求直线斜率为\(-1/2\)：\(y-2=-(1/2)(x-1)\)，即\(x+2y-5=0\)．故选：A．}
\end{ansblock}

\tihao{3}\zu{圆的弦长问题×单选} \tieside{简单(知识点二)}若圆\(x^{2}+y^{2}-2x+4y+m=0\)截直线\(x+y-3=0\)所得弦长为2，则实数\(m\)的值为\nobreak(\kongwei)
\bindopt {\fontsize{10.5pt}{\qpTJgLeadLiti}\selectfont \makebox[\dimexpr0.25\linewidth-0.25\lxhang\relax][l]{A．\(-1\)}\makebox[\dimexpr0.25\linewidth-0.25\lxhang\relax][l]{B．\(-2\)}\makebox[\dimexpr0.25\linewidth-0.25\lxhang\relax][l]{C．\(-4\)}\makebox[\dimexpr0.25\linewidth-0.25\lxhang\relax][l]{D．\(-31\)}\par}

\begin{ansblock}[2章-练-课时08-T3]
% ans:2章-练-课时08-T3
\ansitem{3}{C}
\ansnote{详解}{圆心\((1,-2)\)，\(r^{2}=5-m\)．圆心到直线的距离\(d=|1-2-3|/\sqrt{2}=2\sqrt{2}\)．由半弦长平方\(1=r^{2}-d^{2}\)，得\(1=5-m-8\)，解得\(m=-4\)．故选：C．}
\end{ansblock}

\tihao{4}\zu{圆的弦长问题×单选} \tieside{中档(知识点二)}直线\(l\)：\(y=x+m\)与圆\(x^{2}+y^{2}=4\)相交于\(A\)，\(B\)两点，若\(|AB|\geqslant 2\sqrt{3}\)，则实数\(m\)的取值范围为\nobreak(\kongwei)
\bindopt {\fontsize{10.5pt}{\qpTJgLeadLiti}\selectfont \makebox[\dimexpr0.5\linewidth-0.5\lxhang-0.25em\relax][l]{A．\([-2,2]\)}\makebox[\dimexpr0.5\linewidth-0.5\lxhang-0.25em\relax][l]{B．\([-\sqrt{2},\sqrt{2}]\)}\par}
\bindopt {\fontsize{10.5pt}{\qpTJgLeadLiti}\selectfont \makebox[\dimexpr0.5\linewidth-0.5\lxhang-0.25em\relax][l]{C．\([-1,1]\)}\makebox[\dimexpr0.5\linewidth-0.5\lxhang-0.25em\relax][l]{D．\([-\sqrt{2}/2,\ \sqrt{2}/2]\)}\par}

\begin{ansblock}[2章-练-课时08-T4]
% ans:2章-练-课时08-T4
\ansitem{4}{B}
\ansnote{详解}{\(|AB|\geqslant 2\sqrt{3}\)，则半弦长\(\geqslant\sqrt{3}\)，弦心距\(d\leqslant\sqrt{4-3}=1\)．又\(d=|m|/\sqrt{2}\)，故\(|m|\leqslant\sqrt{2}\)，即\(m\in[-\sqrt{2},\sqrt{2}]\)．故选：B．}
\end{ansblock}

\tihao{5}\zu{直线与圆的位置关系×单选} \tieside{简单(知识点一)}直线 \(l\)：\(y=k(x-1)+1\) 和圆 \(x^{2}+y^{2}-2y=0\) 的位置关系为\nobreak(\kongwei)

\bindopt {\fontsize{10.5pt}{\qpTJgLeadLiti}\selectfont \makebox[\dimexpr0.5\linewidth-0.5\lxhang-0.25em\relax][l]{A．相交}\makebox[\dimexpr0.5\linewidth-0.5\lxhang-0.25em\relax][l]{B．相切或相交}\par}
\bindopt {\fontsize{10.5pt}{\qpTJgLeadLiti}\selectfont \makebox[\dimexpr0.5\linewidth-0.5\lxhang-0.25em\relax][l]{C．相离}\makebox[\dimexpr0.5\linewidth-0.5\lxhang-0.25em\relax][l]{D．相切}\par}

\begin{ansblock}[2章-练-课时08-T5]
% ans:2章-练-课时08-T5
\ansitem{5}{A}
\ansnote{详解}{圆即\(x^{2}+(y-1)^{2}=1\)，圆心\((0,1)\)．直线\(y=k(x-1)+1\)过定点\((1,1)\)（\(x=1\)时\(y=1\)），恰为圆心，故圆心到直线的距离\(d=0<r\)，直线与圆恒相交．故选：A．}
\end{ansblock}

\tihao{6}\zu{圆的切线方程×单选} \tieside{简单(知识点三)}与圆\(x^{2}+y^{2}-4x+3=0\)相切，且在\(x\)、\(y\)轴上截距相等的直线共有\nobreak(\kongwei)
\bindopt {\fontsize{10.5pt}{\qpTJgLeadLiti}\selectfont \makebox[\dimexpr0.25\linewidth-0.25\lxhang\relax][l]{A．1条}\makebox[\dimexpr0.25\linewidth-0.25\lxhang\relax][l]{B．2条}\makebox[\dimexpr0.25\linewidth-0.25\lxhang\relax][l]{C．3条}\makebox[\dimexpr0.25\linewidth-0.25\lxhang\relax][l]{D．4条}\par}

\begin{ansblock}[2章-练-课时08-T6]
% ans:2章-练-课时08-T6
\ansitem{6}{D}
\ansnote{详解}{圆即\((x-2)^{2}+y^{2}=1\)，圆心\((2,0)\)，\(r=1\)．①过原点（两截距同为0）：设\(y=kx\)，\(d=|2k|/\sqrt{1+k^{2}}=1\)，得\(k^{2}=1/3\)，有2条；②不过原点：设\(x+y=a\)，\(d=|2-a|/\sqrt{2}=1\)，得\(a=2\pm\sqrt{2}\)（均非0），有2条．共4条，故选：D．}
\end{ansblock}

\tihao{7}\zu{圆的切线方程×单选} \tieside{简单(知识点三)}若过点\((2,1)\)的圆与两坐标轴都相切，则圆心到直线\(2x-y-3=0\)的距离为\nobreak(\kongwei)
\bindopt {\fontsize{10.5pt}{\qpTJgLeadLiti}\selectfont \makebox[\dimexpr0.25\linewidth-0.25\lxhang\relax][l]{A．\(\sqrt{5}/5\)}\makebox[\dimexpr0.25\linewidth-0.25\lxhang\relax][l]{B．\(2\sqrt{5}/5\)}\makebox[\dimexpr0.25\linewidth-0.25\lxhang\relax][l]{C．\(3\sqrt{5}/5\)}\makebox[\dimexpr0.25\linewidth-0.25\lxhang\relax][l]{D．\(4\sqrt{5}/5\)}\par}

\begin{ansblock}[2章-练-课时08-T7]
% ans:2章-练-课时08-T7
\ansitem{7}{B}
\ansnote{详解}{圆与两轴都相切且过点\((2,1)\)（第一象限），则圆心\((a,a)\)，半径\(a\)（\(a>0\)）．由\((2-a)^{2}+(1-a)^{2}=a^{2}\)得\(a^{2}-6a+5=0\)，\(a=1\)或\(a=5\)．\(a=1\)：圆心\((1,1)\)，\(d=|2-1-3|/\sqrt{5}=2\sqrt{5}/5\)；\(a=5\)：圆心\((5,5)\)，\(d=|10-5-3|/\sqrt{5}=2\sqrt{5}/5\)．两者相同，故\(d=2\sqrt{5}/5\)，故选：B．}
\end{ansblock}

% ============================================================
% 综合提升（题8~14＝T8~T14：单选3（T8~T10）＋填空4（T11~T14））
% ============================================================
\dangbadge{综}{合}{提}{升}

\tihao{8}\zu{直线与圆的位置关系×单选} \tieside{简单(知识点一)}已知集合\(A=\{(x,y)\mid x^{2}+y^{2}=1\}\)，集合\(B=\{(x,y)\mid y=|x|-1\}\)，则集合\(A\cap B\)的真子集的个数为\nobreak(\kongwei)
\bindopt {\fontsize{10.5pt}{\qpTJgLeadLiti}\selectfont \makebox[\dimexpr0.25\linewidth-0.25\lxhang\relax][l]{A．3}\makebox[\dimexpr0.25\linewidth-0.25\lxhang\relax][l]{B．4}\makebox[\dimexpr0.25\linewidth-0.25\lxhang\relax][l]{C．7}\makebox[\dimexpr0.25\linewidth-0.25\lxhang\relax][l]{D．8}\par}

\begin{ansblock}[2章-练-课时08-T8]
% ans:2章-练-课时08-T8
\ansitem{8}{C}
\ansnote{详解}{\(x\geqslant 0\)时，\(y=x-1\)，代入圆方程：\(2x^{2}-2x=0\)，得\(x=0\)（\(y=-1\)）或\(x=1\)（\(y=0\)）；\(x<0\)时，\(y=-x-1\)，代入：\(2x^{2}+2x=0\)，得\(x=-1\)（\(y=0\)）（\(x=0\)不在此段）．交点\((0,-1)\)，\((1,0)\)，\((-1,0)\)共3个，真子集\(2^{3}-1=7\)个．故选：C．}
\end{ansblock}

\tihao{9}\zu{圆的切线方程×单选} \tieside{中档(知识点三)}过直线\(y=2x-3\)上的点作圆\(C\)：\(x^{2}+y^{2}-4x+6y+12=0\)的切线，则切线长的最小值为\nobreak(\kongwei)
\bindopt {\fontsize{10.5pt}{\qpTJgLeadLiti}\selectfont \makebox[\dimexpr0.25\linewidth-0.25\lxhang\relax][l]{A．\(\sqrt{55}/5\)}\makebox[\dimexpr0.25\linewidth-0.25\lxhang\relax][l]{B．\(\sqrt{19}\)}\makebox[\dimexpr0.25\linewidth-0.25\lxhang\relax][l]{C．\(2\sqrt{5}\)}\makebox[\dimexpr0.25\linewidth-0.25\lxhang\relax][l]{D．\(\sqrt{21}\)}\par}

\begin{ansblock}[2章-练-课时08-T9]
% ans:2章-练-课时08-T9
\ansitem{9}{A}
\ansnote{详解}{圆\(C\)即\((x-2)^{2}+(y+3)^{2}=1\)，圆心\(C(2,-3)\)，\(r=1\)．切线长\(=\sqrt{|PC|^{2}-r^{2}}\)，最小当且仅当\(|PC|\)最小，即\(C\)到直线的距离\(|4+3-3|/\sqrt{5}=4/\sqrt{5}\)．最小切线长\(=\sqrt{16/5-1}=\sqrt{11/5}=\sqrt{55}/5\)．故选：A．}
\end{ansblock}

\tihao{10}\zu{圆的切线方程×单选} \tieside{简单(知识点三)}已知直线\(l\)：\(x+ay-1=0\)是圆\(C\)：\(x^{2}+y^{2}-6x-2y+1=0\)的对称轴，过点\(A(-1,a)\)作圆\(C\)的一条切线，切点为\(B\)，则\(|AB|=\)\nobreak(\kongwei)
\bindopt {\fontsize{10.5pt}{\qpTJgLeadLiti}\selectfont \makebox[\dimexpr0.25\linewidth-0.25\lxhang\relax][l]{A．\(1\)}\makebox[\dimexpr0.25\linewidth-0.25\lxhang\relax][l]{B．\(2\)}\makebox[\dimexpr0.25\linewidth-0.25\lxhang\relax][l]{C．\(4\)}\makebox[\dimexpr0.25\linewidth-0.25\lxhang\relax][l]{D．\(8\)}\par}

\begin{ansblock}[2章-练-课时08-T10]
% ans:2章-练-课时08-T10
\ansitem{10}{C}
\ansnote{详解}{圆\(C\)即\((x-3)^{2}+(y-1)^{2}=9\)，圆心\((3,1)\)，\(r=3\)．对称轴过圆心：\(3+a-1=0\)，得\(a=-2\)，故\(A(-1,-2)\)．\(|AC|=\sqrt{16+9}=5\)，切线长\(|AB|=\sqrt{25-9}=4\)．故选：C．}
\end{ansblock}

\tihao{11}\zu{直线与圆的位置关系×填空} \tieside{简单(知识点一)}直线\(x\cos\theta+y\sin\theta=r\)和\(x^{2}+y^{2}=r^{2}\)的位置关系是\kongda{相切}．

\begin{ansblock}[2章-练-课时08-T11]
% ans:2章-练-课时08-T11
\ansitem{11}{相切}
\ansnote{详解}{圆心\((0,0)\)，半径\(r\)．圆心到直线的距离\(d=|-r|/\sqrt{\cos^{2}\theta+\sin^{2}\theta}=r\)，恒等于半径，故直线与圆恒相切．}
\end{ansblock}

\tihao{12}\zu{直线与圆的位置关系×填空} \tieside{中档(知识点一)}若圆\(C\)：\(x^{2}+y^{2}-2mx+2\sqrt{m}\,y+2=0\)与\(x\)轴有公共点，则实数\(m\)的取值范围是\kongda{[√2, +∞)}．

\begin{ansblock}[2章-练-课时08-T12]
% ans:2章-练-课时08-T12
\ansitem{12}{[√2, +∞)}
\ansnote{详解}{表为圆需\(\sqrt{m}\)有意义，故\(m\geqslant 0\)，且\(r^{2}=m^{2}+m-2>0\)，在\(m\geqslant 0\)下即\(m>1\)．与\(x\)轴有公共点\(\Longleftrightarrow|\sqrt{m}|\leqslant r\)，即\(m\leqslant m^{2}+m-2\)，得\(m^{2}\geqslant 2\)，故\(m\geqslant\sqrt{2}\)．所求取值范围为\([\sqrt{2},+\infty)\)．}
\end{ansblock}

\tihao{13}\zu{直线与圆的位置关系×填空} \tieside{简单(知识点一)}已知圆\(C\)：\((x-1)^{2}+y^{2}=1\)，点\(A(-2,0)\)及点\(B(3,a)\)，从点\(A\)处观察点\(B\)，要使视线不被圆\(C\)挡住，则实数\(a\)的取值范围为\kongda{(−∞, −5√2/4)∪(5√2/4, +∞)}.

\begin{ansblock}[2章-练-课时08-T13]
% ans:2章-练-课时08-T13
\ansitem{13}{(−∞, −5√2/4)∪(5√2/4, +∞)}
\ansnote{详解}{视线不被挡住\(\Longleftrightarrow\)直线\(AB\)与圆\(C\)相离．\(k_{AB}=a/5\)，直线\(AB\)：\(ax-5y+2a=0\)．圆心\((1,0)\)到直线\(AB\)的距离\(d=|3a|/\sqrt{a^{2}+25}>1\)，即\(9a^{2}>a^{2}+25\)，\(8a^{2}>25\)，\(|a|>5/(2\sqrt{2})=5\sqrt{2}/4\)．故\(a\)的取值范围为\((-\infty,-5\sqrt{2}/4)\cup(5\sqrt{2}/4,+\infty)\)．}
\end{ansblock}

\tihao{14}\zu{圆的切线方程×填空} \tieside{简单(知识点三)}若斜率为\(\sqrt{3}\)的直线与\(y\)轴交于点\(A\)，与圆\(x^{2}+(y-1)^{2}=1\)相切于点\(B\)，则\(|AB|=\)\kongda{√3}．

\begin{ansblock}[2章-练-课时08-T14]
% ans:2章-练-课时08-T14
\ansitem{14}{√3}
\ansnote{详解}{设直线\(y=\sqrt{3}x+b\)，则\(A(0,b)\)．相切时圆心\((0,1)\)到直线的距离\(|b-1|/2=1\)，得\(b=3\)或\(b=-1\)，\(|AC|=2\)（\(C(0,1)\)）．切线长\(|AB|=\sqrt{|AC|^{2}-r^{2}}=\sqrt{4-1}=\sqrt{3}\)．}
\end{ansblock}

% ============================================================
% 思维探索（题15~16＝T15~T16：解答2，居末档照库序）
% ============================================================
\dangbadge{思}{维}{探}{索}

\tihao{15}\zu{圆的切线方程＋弦长综合×解答} \tieside{中档(知识点三)}已知点\(M(3,1)\)，直线\(ax-y+4=0\)及圆\(C\)：\((x-1)^{2}+(y-2)^{2}=4\)．
(1) 求过点\(M\)的圆\(C\)的切线方程；
(2) 若直线\(ax-y+4=0\)与圆\(C\)相切，求实数\(a\)的值；
(3) 若直线\(ax-y+4=0\)与圆\(C\)相交于\(A\)、\(B\)两点，且弦\(AB\)的长为\(2\sqrt{3}\)，求\(a\)的值．
\liubai[32mm]{解答书写区}

\begin{ansblock}[2章-练-课时08-T15]
% ans:2章-练-课时08-T15
\ansitem{15}{(1) x=3 或 3x−4y−5=0；(2) a=0 或 a=4/3；(3) a=−3/4}
\ansnote{详解}{\(C(1,2)\)，\(r=2\)；\(|MC|=\sqrt{4+1}=\sqrt{5}>2\)，\(M\)在圆外．\par\noindent (1) 斜率不存在时\(x=3\)：圆心到直线的距离\(|3-1|=2=r\)，是切线．斜率存在时设\(y-1=k(x-3)\)，即\(kx-y-3k+1=0\)：\(d=|2k+1|/\sqrt{k^{2}+1}=2\)，得\(4k^{2}+4k+1=4k^{2}+4\)，\(k=3/4\)，切线为\(3x-4y-5=0\)．故切线方程为\(x=3\)或\(3x-4y-5=0\)．\par\noindent (2) \(d=|a+2|/\sqrt{a^{2}+1}=2\)，得\(a^{2}+4a+4=4a^{2}+4\)，即\(3a^{2}-4a=0\)，解得\(a=0\)或\(a=4/3\)．\par\noindent (3) 弦长\(2\sqrt{3}\)，则半弦长\(\sqrt{3}\)，\(d=\sqrt{4-3}=1\)：\(|a+2|/\sqrt{a^{2}+1}=1\)，得\(a^{2}+4a+4=a^{2}+1\)，解得\(a=-3/4\)．}
\end{ansblock}

\tihao{16}\zu{阿波罗尼斯圆问题×解答} \tieside{中档(知识点一)}为了保证我国东海油气田海域的海上平台的生产安全，海事部门在某平台\(O\)的正东方向设立了两个观测站\(A\)、\(B\)（点\(A\)在点\(O\)、点\(B\)之间），它们到平台\(O\)的距离分别为3海里和12海里，记海平面上到两观测站距离\(PA\)、\(PB\)之比为\(1/2\)的点\(P\)的轨迹为曲线\(E\)，规定曲线\(E\)及其内部区域为安全预警区．
(1) 以\(O\)为坐标原点，\(AB\)所在直线为\(x\)轴建立平面直角坐标系，求曲线\(E\)的方程；
(2) 某日在观测站\(B\)处发现，在该海上平台正南\(2\sqrt{11}\)海里的\(C\)处，有一艘轮船正以每小时10海里的速度向北偏东\(30^\circ\)方向航行，如果航向不变，该轮船是否会进入安全预警区？如果不进入，说明理由；如果进入，则它在安全预警区中的航行时间是几小时．
\liubai[24mm]{解答书写区}

\begin{ansblock}[2章-练-课时08-T16]
% ans:2章-练-课时08-T16
\ansitem{16}{(1) x²+y²=36；(2) 一定进入，航行时间 1 小时}
\ansnote{详解}{(1) 以\(O\)为原点，则\(A(3,0)\)，\(B(12,0)\)．由\(|PA|/|PB|=1/2\)得\(4[(x-3)^{2}+y^{2}]=(x-12)^{2}+y^{2}\)，即\(4x^{2}-24x+36+4y^{2}=x^{2}-24x+144+y^{2}\)，整理得\(3x^{2}+3y^{2}=108\)，故曲线\(E\)的方程为\(x^{2}+y^{2}=36\)．\par\noindent (2) \(C(0,-2\sqrt{11})\)．北偏东\(30^\circ\)方向即航线倾角\(60^\circ\)，斜率\(\sqrt{3}\)，航线：\(y+2\sqrt{11}=\sqrt{3}x\)，即\(\sqrt{3}x-y-2\sqrt{11}=0\)．圆心\(O\)到航线的距离\(d=|2\sqrt{11}|/\sqrt{3+1}=\sqrt{11}<6\)，直线与圆相交，故轮船一定进入安全预警区．弦长\(=2\sqrt{36-11}=10\)，航行时间\(t=10/10=1\)（小时）．}
\end{ansblock}

% ---- 尾块（\tailfill 置 \end{multicols} 前；无参＝内容尾随框，每件恰 1 处） ----
\tailfill

\end{multicols}
\end{document}
